%
%============= CHAPTER BREAK =================%
%
\chapter{Herodian Hist. (244--251 CE)}
%
\vspace*{3mm}
%
\justifying

\noindent{}As previously noted, Herodian the Historian (early 3rd century CE), %
should not be confused with Aelius Herodian the Alexandrian grammarian %
and son of the grammarian Apollonius Dyscolus. %
The historian claimed to have lived through some 60 or 70 years of events %
which he narrated in his most famous work, %
entitled \textit{From the Departure of Divine Marcus [Aurelius]} / %
\textit{Ab excessu divi Marci}, %
written under Philip the Arab (244--249 CE) and/or Decius (249--251 CE).\endnote{%
Baumann et al 1025:1n2, with further literature.} \\[1.5mm] %

\noindent{}\textbf{διαγρηγορέω} / \textit{diagrēgoreō} / “to keep awake|vigil” appears in Luke 9:32, %
apparently referring to “Peter” and “those with him” observing a vigil during the transfiguration of Jesus. %
The earliest non-Christian usage of the term is found in Herodian the Historian, %
in a scene where the term describes a harrowing civil war battle %
between Septimius Severus and Pescennius Niger in 193--194 CE. %

\vspace*{3mm}
 \begin{quote}συνέβη δὲ καὶ τῶν στρατοπέδων τοῦ Σεβήρου τε καὶ %
 <τοῦ> Νίγρου μὴ τὴν σύνοδον μόνον κατ’ ἐκεῖνο γενέσθαι τὸ χωρίον, %
 ἀλλὰ καὶ τὴν τύχην ὁμοίαν τῆς μάχης. %
 ἀντιστρατοπεδευσάμενοι γὰρ ἑκατέρωθεν περὶ ἑσπέραν, %
 πάσης τῆς νυκτὸς ἐν φροντίσιν ἑκάτεροι καὶ δέει \textbf{διαγρηγορήσαντες}, %
 ἅμα ἡλίῳ ἀνίσχοντι ἐπ’ ἀλλήλους ἠπείγοντο, %
 παρορμώντων ἑκατέρωθεν τῶν στρατηγῶν. %
 προθυμίᾳ δὴ πάσῃ ἐνέπιπτον ὡς ὑπὲρ λοιπῆς καὶ τελευταίας ἐκείνης μάχης, %
 κἀκεῖ τῆς τύχης διακρινούσης τὸν βασιλέα. \\[0.1mm] %
 Now it came about that not only did the assembly of the armies of both Severus %
 and of Niger happen in that very place, but also the same fate of the battle. %
 For after encamping opposite [each other] on each side for the evening, %
 each of them at attention the whole night and in dread \textbf{keeping vigil}, %
 all at once when the sun rose they rushed upon each other, %
 incited from both sides by the commanders. %
 With all intensity they were falling in, as for that last and conclusive battle, %
 eʼen there where Fortune would determine the king.\end{quote} % 
\vspace*{3mm}

\noindent{}In Luke 9:32, we have yet another example of an unusual term and %
apparent neologism not present in the substantial parallels in Marcion's \textit{Evangelion}, %
canonical Mark, and canonical Matthew. %

\vspace*{3mm} 
 \begin{quote}ὁ δὲ Πέτρος καὶ οἱ σὺν αὐτῷ ἦσαν βεβαρημένοι ὕπνῳ· %
 \textbf{διαγρηγορήσαντες} δὲ εἶδον τὴν δόξαν αὐτοῦ καὶ τοὺς δύο ἄνδρας τοὺς συνεστῶτας αὐτῷ. \\[1mm] %
 Now Peter and those with him had been weighed down with sleep. %
 But \textbf{after keeping vigil}, they saw the glory of his and the two men who stood alongside him. %
 \end{quote} % 
\vspace*{2mm}

If aware of the carefully crafted imagery of Herodian the Historian, %
then some mid-3rd century (!) scribal editor of canonical Luke used it %
to depict Peter and companions not only as heroic in religious devotion, %
but also as vivid rivals of the vigilant militancy of a recent Roman civil war. %
If this reading is valid, then Jesus becomes the bringer of peace amidst ancient military leaders like Moses, %
lethal prophets like Elijah, and his own militarily overeager yet inept Galilean followers. %

There is a variant here to consider, though, %
since 𝔓\textsuperscript{45} has διεγρηγορήσαντες / \textit{diegrēgorēsantes} from διεγείρω / \textit{diegeirō} / “to wake up”, %
in a perfect form similar to that attested by Aristophanes Gram. (3rd--2nd c. BCE): %
διεγρηγόρασι / \textit{diegrēgorasi} (\textit{hist.anim.} 2.306). %
Otherwise, manuscripts of Luke 9:32 are uniform in using the lemma found first in Herodian the Historian. %

